\section{随机变量}
\begin{definition}[随机变量]
设 $S = \{e\}$ 是试验的样本空间，如果量 $X$ 是定义在 $S$ 上的一个单值实值函数，则称 $X$ 为\textbf{随机变量}（r.v. $X$）。
\begin{itemize}
	\item \textbf{映射}：$X: S \to \mathbb{R}$。
	\item \textbf{引入意义}：将随机试验的结果数量化，描述随机事件。
	\item \textbf{分类}：离散型随机变量、非离散型（连续型、奇异型/混合型）。
\end{itemize}
\end{definition}

\section{离散型随机变量的分布律}
\begin{definition}[分布律]
若 $X$ 取值 $x_1, x_2, \ldots, x_k, \ldots$ 且取这些值的概率依次为 $p_1, p_2, \ldots, p_k, \ldots$，则称
\begin{equation}
	P\{X = x_k\} = p_k \quad (k = 1, 2, \ldots)
	\label{eq:第二章:分布律}
\end{equation}
为 $X$ 的\textbf{分布律}。
\end{definition}

\begin{property}[分布律的性质]
\begin{enumerate}
	\item $p_k \ge 0, \ k = 1, 2, 3, \dots$;
	\item $\sum_{k \ge 1} p_k = 1$.
\end{enumerate}
\end{property}

\subsection{常见的离散型随机变量的分布律}
\begin{definition}[两点分布]
\begin{center}
\begin{tabular}{c|cc}
\toprule
$X$ & 1 & 0 \\
\midrule
$p_k$ & $1-p$ & $p$ \\
\bottomrule
\end{tabular}
\end{center}
\end{definition}

\begin{definition}[二项分布]
记作 $X \sim b(n, p)$。其中：
\begin{itemize}
	\item $X$ 表示 $n$ 重贝努里试验中事件 $A$ 发生的次数；
	\item $n$ 表示伯努利试验发生的次数；
	\item $p = P(A)$。
\end{itemize}
\begin{equation}
	P\{X = k\} = C_n^k\, p^k (1 - p)^{n - k}, \quad (k = 0, 1, 2, 3, \dots, n)
	\label{eq:第二章:二项分布}
\end{equation}
\end{definition}

\begin{definition}[泊松分布]
泊松分布是 $n$ 很大、$p$ 很小的二项分布的极限。如果一个随机变量 $X$ 服从参数为 $\lambda$ 的泊松分布，即 $X \sim \pi(\lambda)$，其概率质量函数为：
\begin{equation}
	P(X = k) = \frac{e^{-\lambda} \lambda^k}{k!}, \quad (k = 0, 1, 2, \dots)
	\label{eq:第二章:泊松分布}
\end{equation}
\begin{itemize}
	\item $\lambda$（Lambda）：这是唯一的参数，代表事件发生的\textbf{平均速率}（或平均发生次数）。
	\item $e$：自然常数（约等于 2.718）。
	\item $k$：你想要计算的实际发生次数。
\end{itemize}
\end{definition}

\begin{example}[泊松分布：问讯次数]
$X \sim P(\lambda)$，已知 $P\{X=0\} = e^{-6}$。
\begin{itemize}
	\item \textbf{求参数 $\lambda$}：$P\{X=0\} = \frac{\lambda^{0}}{0!} e^{-\lambda} = e^{-\lambda}$，故 $\lambda = 6$。
	\item \textbf{求 $P\{X \ge 2\}$}：
	\begin{itemize}
		\item \textbf{泊松分布公式}：$P\{X=k\} = \frac{\lambda^{k}}{k!} e^{-\lambda}$
		\begin{equation}
			P\{X \ge 2\} = 1 - P\{X < 2\} = 1 - \big(P\{X=0\} + P\{X=1\}\big)
			\label{eq:第二章:泊松例1}
		\end{equation}
		\begin{equation}
			P\{X \ge 2\} = 1 - \left(e^{-6} + \frac{6^{1}}{1!} e^{-6}\right) = 1 - 7 e^{-6} \approx 0.9826
			\label{eq:第二章:泊松例2}
		\end{equation}
	\end{itemize}
\end{itemize}
\end{example}

\section{随机变量的分布函数}
\begin{definition}[分布函数]
设 $X$ 是一个随机变量，\textbf{分布函数} $F(x)$ 定义为：
\begin{equation}
	F(x) = P\{X \le x\}, \quad -\infty < x < +\infty
	\label{eq:第二章:分布函数}
\end{equation}
\end{definition}

\begin{property}[分布函数的性质]
\begin{enumerate}
	\item 单调不减性（Monotonicity）：若 $x_1 < x_2$，则 $F(x_1) \le F(x_2)$；
	\item 归一性（Normalization / Boundedness）：$0 \le F(x) \le 1$，且 $F(-\infty) = 0, F(+\infty) = 1$；
	\item 右连续性（Right-continuity）：$\lim_{x \to x_0^+} F(x) = F(x_0)$。
\end{enumerate}
\end{property}

\begin{keybox}[分布函数常用公式]
\begin{itemize}
	\item $P(a < X \le b) = F(b) - F(a)$
	\item $P(X > a) = 1 - F(a)$
	\item $P(X = a) = F(a) - F(a-0)$
	\item $P(X < a) = F(a-0)$
\end{itemize}
\end{keybox}

\section{连续型随机变量及其概率密度}
\begin{definition}[概率密度函数]
若 $X$ 的分布函数为 $F(x)$，存在非负函数 $f(x)$，使得 $F(x) = \int_{-\infty}^{x} f(t)\,dt$，则称 $f(x)$ 为 $X$ 的\textbf{概率密度函数}。
\begin{enumerate}
	\item $f(x) \ge 0$；
	\item 归一性：$\int_{-\infty}^{\infty} f(x)\,dx = 1$。
\end{enumerate}
概率计算：$P\{a < X < b\} = \int_{a}^{b} f(x)\,dx = F(b) - F(a)$。
\end{definition}

\subsection{几个常用的连续型分布}
\begin{definition}[均匀分布]
如果随机变量 $X$ 在区间 $[a, b]$ 上服从均匀分布，记作 $X \sim U(a, b)$。它的概率密度函数是一条水平直线：
\begin{equation}
	f(x) = \begin{cases} \dfrac{1}{b - a}, & a \le x \le b \\[1ex] 0, & \text{其他} \end{cases}
	\label{eq:第二章:均匀密度}
\end{equation}
\begin{equation}
	F(x) = P(X \le x) = \begin{cases} 0, & x < a \\[1ex] \dfrac{x - a}{b - a}, & a \le x < b \\[1ex] 1, & x \ge b \end{cases}
	\label{eq:第二章:均匀分布函数}
\end{equation}
\end{definition}

\begin{definition}[指数分布]
如果随机变量 $X$ 服从参数为 $\lambda$ 的指数分布（其中 $\lambda > 0$），记作 $X \sim E(\lambda)$。它的概率密度函数为：
\begin{equation}
	f(x) = \begin{cases} \lambda e^{-\lambda x}, & x \ge 0 \\[1ex] 0, & x < 0 \end{cases}
	\label{eq:第二章:指数密度}
\end{equation}
\begin{equation}
	F(x) = P(X \le x) = \begin{cases} 1 - e^{-\lambda x}, & x \ge 0 \\[1ex] 0, & x < 0 \end{cases}
	\label{eq:第二章:指数分布函数}
\end{equation}
\end{definition}

\begin{definition}[正态分布]
如果随机变量 $X$ 服从均值为 $\mu$、方差为 $\sigma^2$ 的正态分布，记作 $X \sim N(\mu, \sigma^2)$。其密度函数为：
\begin{equation}
	f(x) = \frac{1}{\sqrt{2\pi}\,\sigma} e^{-\frac{(x-\mu)^2}{2\sigma^2}}
	\label{eq:第二章:正态密度}
\end{equation}
当 $\mu = 0, \sigma = 1$ 时，称为标准正态分布，记为 $Z \sim N(0, 1)$。
\begin{equation}
	\varphi(x) = \frac{1}{\sqrt{2\pi}} e^{-\frac{x^{2}}{2}}, \quad -\infty < x < +\infty
	\label{eq:第二章:标准正态密度}
\end{equation}
\begin{equation}
	\Phi(x) = P\{X \le x\} = \int_{-\infty}^{x} \varphi(t)\,dt
	\label{eq:第二章:标准正态分布函数}
\end{equation}
其分布函数通常用专用符号 $\Phi(x)$ 表示。任何一般的正态分布都可以通过以下方式标准化：
\begin{equation}
	Z = \frac{X - \mu}{\sigma}
	\label{eq:第二章:标准化}
\end{equation}
\end{definition}

\begin{example}[求参数 $A$、概率、分布函数]
已知 $f(x) = \begin{cases} A x^{2}, & 0 < x < 1 \\ 0, & \text{其他} \end{cases}$
\begin{enumerate}
	\item \textbf{求参数 $A$}：由 $\int_{0}^{1} A x^{2} dx = 1 \Rightarrow A \cdot \dfrac{x^{3}}{3}\Big|_{0}^{1} = \dfrac{A}{3} = 1 \Rightarrow A = 3$。
	\item \textbf{求 $P\{0.5 < X < 3\}$}：$P\{0.5 < X < 3\} = \int_{0.5}^{1} 3 x^{2} dx = x^{3}\Big|_{0.5}^{1} = 1 - (0.5)^{3} = 0.875$。
	\item \textbf{求分布函数 $F(x)$}：
	\begin{equation}
		F(x) = \begin{cases} 0, & x \le 0 \\[1ex] \displaystyle\int_{0}^{x} 3 t^{2} dt = x^{3}, & 0 < x < 1 \\[1ex] 1, & x \ge 1 \end{cases}
		\label{eq:第二章:例1分布函数}
	\end{equation}
\end{enumerate}
\end{example}

\subsection{无记忆性}
\begin{definition}[无记忆性]
如果随机变量 $X$ 具有无记忆性，对于任何正数 $s$ 和 $t$，以下等式成立：
\begin{equation}
	P(X > s + t \mid X > s) = P(X > t)
	\label{eq:第二章:无记忆性}
\end{equation}
\end{definition}

\section{随机变量函数的分布}
\begin{theorem}[随机变量函数的分布：公式法]
若 $X \sim f_X(x)$，$Y = g(X)$ 是严格单调可导函数，则 $Y$ 的概率密度函数 $f_Y(y)$ 为：
\begin{equation}
	f_Y(y) = f_X\big[g^{-1}(y)\big] \cdot \left|\frac{d\,g^{-1}(y)}{d y}\right|
	\label{eq:第二章:公式法}
\end{equation}
\end{theorem}

\begin{example}[线性变换]
已知 $X \sim N(\mu, \sigma^{2})$，求 $Y = \frac{X-\mu}{\sigma}$ 的概率密度。
\begin{itemize}
	\item \textbf{反函数}：$x = g^{-1}(y) = \sigma y + \mu$。
	\item \textbf{导数}：$\left|\dfrac{d g^{-1}(y)}{d y}\right| = \sigma$。
	\item \textbf{结果}：$f_Y(y) = f_X(\sigma y + \mu) \cdot \sigma = \dfrac{1}{\sqrt{2\pi}} e^{-\frac{y^{2}}{2}}$。
	\item \textbf{结论}：$Y \sim N(0, 1)$。
\end{itemize}
\end{example}
